Der Zusammenhang mit 1. wird hergestellt durch die Bemerkung, daß die Potenzsummen $S_1,\ldots,S_h$ sich durch die elementaren symmetrischen Funktionen von $\xi_1,\ldots,\xi_h$ ersetzen lassen, was auf das dort gegebene volle System der $G_{\alpha\alpha_1\ldots\alpha_n}(x)$ führt. Beide Resultate zeigen, daß sich alle Invarianten ganz und rational durch diejenigen ausdrücken lassen, deren Grad in den $x$ die Ordnung $h$ der Gruppe nicht übersteigt.
